\documentclass[twoside]{article}
\usepackage{mmpr}

\begin{document}

\Russian
\title{Шаблон оформления тезисов доклада на~конференцию ИОИ-2026}
\author{Автор~И.\,О.}{Автор Имя Отчество$^{1, 2}$}{author\_email@site.ru}
\author{Соавтор~И.\,О.}{Соавтор Имя Отчество$^2$\speaker}{coauthor@site.ru}
\organization{%
    $^1$Город, Организация\par
    $^2$Город, Организация}
\maketitle

Сборник тезисов докладов издается к~началу конференции в печатном виде.
После завершения конференции сборник регистрируется в РИНЦ, и~публикуется его электронная версия.

Тезисы доклада принимаются на~одном языке: русском или английском, и должны занимать до~трех страниц 
 (приблизительно 6000 символов).
Программный комитет оставляет за~собой право сократить текст тезисов,
если он не~умещается на~трех страницах.

Все поля шаблона должны быть заполнены.
Фамилия с~инициалами нужна для авторского указателя.
Фамилия с~полными именем и~отчетсвом~--- для подписи статьи.
Обязательно указание электронных адресов авторов.
Номер сноски для указания аффилированных организаций ставится после отчества без пробела.
Командой \verb|\speaker| необходимо отметить того автора, чье выступление с~докладом планируется.

Тезисы доклада могут содержать формулы, таблицы, иллюстрации.
Тезисы доклада не должны содержать разделов, списков, сносок, плавающих иллюстраций.
Ссылку на~источники финансирования можно указать в~последнем абзаце.
Список литературы может содержать ссылки по~теме доклада.
Для электронных журналов URL статьи указывается обязательно.

На~конференции может быть представлен абсолютно новый или уже опубликованный свежий результат.
В последнем случае текст тезисов обязательно должен содержать ссылку на~опубликованную статью.
Решение о~принятии или отклонении доклада на~конференцию ИОИ-2026
будет приниматься Программным комитетом на~основании текста тезисов.
По~результатам выступления на~конференции Программный комитет может рекомендовать доклад для публикации в~виде расширенной статьи.
В~случае, если доклад делается по~уже опубликованному результату, такая статья должна существенно отличаться от~предыдущей.

Работа поддержана грантом РНФ \No\,00-00-00000.

\begin{thebibliography}{1}
\bibitem{rus_article_rus}
    \emph{Автор\;И.\,О.}
    Название статьи~//
    Название журнала.~--- 2000.~--- Т.\,1.,\,\No\,1.~--- С.\,5--25.
\bibitem{rus_article_eng}
    \emph{Author\;N.}
    Article Title~//
    Journal Title.~--- 2000.~--- Vol.\,1, No.\,1.~--- Pp.\,5--25.
\bibitem{rus_book_rus}
    \emph{Автор\;И.\,О.}
    Название монографии.~--- Город: Издательство, 2000.~--- 25~с.
\bibitem{rus_book_eng}
    \emph{Author\;N.}
    Book Title.~--- City: Publisher, 2000.~--- 25~p.
\end{thebibliography}

\English
\title{Abstract Template for the  IDP-2026 Conference}
\author{Author~N.}{Author Name$^{1, 2}$}{author\_email@site.ru}
\author{Coauthor~N.}{Coauthor Name$^2$\speaker}{coauthor@site.ru}
\organization{%
    $^1$City, Institution\par
    $^2$City, Institution}
\maketitle

The book of abstracts is published by the beginning of the conference in printed form.
After the conference it is registered in the Russian Science Citation Index, and its electronic version is released.

Abstracts of reports are accepted in either Russian or English, and should be up to three pages long (approximately 6000 characters).
The program committee reserves the right to shorten the text of the abstract if it does not fit the limit.

All fields of the template must be filled in.
The last name with initials is required for the author index.
The full name is needed for the article title.
The e-mail addresses of the authors are also mandatory.
The footnote number for affiliated organizations should be placed after the name without a space.
The \verb|\speaker| command must be used to mark the presenting author.

The abstract of the report may contain formulas, tables, illustrations.
It should not contain sections, lists, footnotes, floating illustrations.
the last paragraph may contain a description of funding sources.
The abstract may contain references on the topic of the report.
For electronic journals, the URL of the article is mandatory.

The conference accepts completely new or already published fresh results.
In the latter case, the text of the abstract must necessarily contain a link to the published article.
The decision on accepting or rejecting a report for the IDP-2026 conference
will be made by the Program Committee based on the text of the abstract.
Based on the results of the presentation at the conference, the Program Committee can recommend the report for publication as an extended article.
If the report is made on an already published result, such an article must differ significantly from the previous one.

The research is funded by RSF, project No.\,00-00-00000.

\begin{thebibliography}{1}
\bibitem{eng_article_rus}
    \emph{Avtor\;I.\,O.}
    Nazvanie stat'i~//
    Nazvanie zhurnala.~--- 2000.~--- Vol.\,1.,\,No.\,1.~--- Pp.\,5--25. (In Russian)
\bibitem{eng_article_eng}
    \emph{Author\;N.}
    Article Title~//
    Journal Title.~--- 2000.~--- Vol.\,1, No.\,1.~--- Pp.\,5--25.
\bibitem{eng_book_rus}
    \emph{Avtor\;I.\,O.}
    Nazvanie monografii.~--- City: Izdatelstvo, 2000.~--- 25~p. (In Russian)
\bibitem{eng_book_eng}
    \emph{Author\;N.}
    Book Title.~--- City: Publisher, 2000.~--- 25~p.
\end{thebibliography}

\end{document} 